% %
% % Alex: I think a good use of this section is to discuss work that
% % isn't already mentioned more naturally elsewhere in the paper.  So stuff like other approaches to fairness and DLPs are already covered by the other sections.  We cite Corbett-Davies a ton.  Kleinberg is cited in connection to notions of fairness and trade-offs between them.  
% %

% \subsection{Disparate Learning Processes}
% calders 


% \subsection{Related Algorithms}
% goh? or is it DLP?  (Alex: I don't think it is)
% bechavod

% \subsection{Other Approaches to Fairness}
% moritz
% or does bechavod go here?
% cynthia

% \alex{With the exception of Cynthia (individual fairness), a survey of other notions already appears in Section 5.3.}

% \subsection{Critical Work}
% nissenbaum
% kleinberg
% corbett

% \begin{itemize}
% \item Friedman Nissenbaum \cite{friedman1996bias}
% \item \textbf{DLPs}
% \item \cite{pedreshi2008discrimination}
% \item \cite{kamiran2009classifying}
% \item \cite{calders2010three}
% \item \cite{luong2011k} %not disp impact
% \item \cite{dwork2012fairness} % individual fairness
% \item \cite{zemel2013learning}
% \item \cite{zafar2017fairness}
% \item \textbf{MITIGATING DISPARATE IMPACT BY USING PROTECTED FEATURE}
% \item \cite{hardt2016equality} % not disparate impact
% \item \cite{dwork2017decoupled} describes decoupling % already cited
% \item \textbf{BACKGROUND ON DISCRIMINATION}
% \item \cite{barocas2016big} % already cited
% \item \cite{USConst} % already cited
% \item \cite{grimmelmann2017incomprehensible} (Zootopia paper) % already cited
% \item \textbf{ANALYSIS}
% \item \cite{kleinberg2016inherent} % already cited
% \item \cite{corbett2017algorithmic} % already cited
% \end{itemize}


% In this section we provide a brief overview 
% of some of the other approaches 
% that have been put forth for trading off 
% between classification performance and impact disparity.  
% One common approach consists of preprocessing 
% or ``massaging'' the training data 
% to reduce the dependence between 
% the resulting model predictions 
% and the sensitive attribute \citep{kamiran2009classifying,kamiran2012data,
% feldman2015certifying,adler2016auditing,
% johndrow2017algorithm}. 
% These methods differ both in terms of 
% what variables are affected by the data processing, 
% and the degree of independence that is achieved.  
% For instance, \citet{kamiran2009classifying} 
% propose flipping the negative labels 
% of some observations in the disadvantaged class.
% \citet{zemel2013learning} proposes 
% learning representations---in this case, cluster assignments---of each example
% such that each example maps to a cluster with some probability, seeking parity in the proportion of each group assigned to each cluster.
% \citet{feldman2015certifying} also investigates 
% transformations of the features $X$ 
% into a new set of features that are 
% constructed to be marginally independent from $Z$.  
% \citet{johndrow2017algorithm} demonstrate 
% how to construct transformations 
% to ensure that the derived features 
% are jointly independent of $Z$, 
% and show that this produces distributional parity 
% of the resulting fitted model. 

% A second common approach 
% is to modify existing classification methods 
% either through post-hoc corrections 
% or in the training stage 
% to constrain the level of impact parity 
% in the resulting model. 
% \citet{kamishima2011fairness, goh2016satisfying, calders2010three,kamiran2010discrimination} 
% consider modifications to methods such as SVM, 
% logistic regression, Naive Bayes, and decision trees.  
% \citet{agarwal2017reductions} 
% show how impact parity constraints 
% can be framed as a cost-sensitive classification problem. 
